% Recuperación de 13_accion_interaccion_absorbedor y83_operador_absorbedor.
% Desarrollo del canal cronológico desde el transporte de rutas.
% Procedencia y condiciones en PROCEDENCIA_PROPAGACION_BIDIRECCIONAL.md.
% Los originales y la sección80 permanecen íntegros. Sin contenido de radión.

\section{Propagación bidireccional y condición de absorción en la frontera}
\label{vii:sec:propagacion-bidireccional}

El transporte APP--TRIT--TPK conserva la orientación de la ruta y la memoria
que distingue su ida de su retorno. Sobre ese antecedente se construyen tres
operaciones sucesivas: la acción de una historia enriquecida, la propagación
de una fuente a través de los transportes realizados y la imposición de una
condición común en la frontera. La inscripción del registro en un espacio de
aparato, establecida en la sección~\ref{vii:sec:observador-aparato}, permite
expresar esas operaciones mediante aplicaciones lineales sin confundir la
representación observable con el estado que conserva sus antecedentes.

La dirección temporal, la conjugación de hojas y la reflexión espacial
tienen dominios distintos. Sus relaciones se fijan primero sobre el registro;
después se especifican los entrelazamientos que utiliza una realización
dinámica. Este orden determina qué información permanece en los propagadores
y qué condición introduce la elección de una frontera absorbente.

\subsection{Acción de rutas e involuciones del registro}
\label{vii:subsec:accion-bidireccional}

Una ruta finita contiene sus posiciones, direcciones cardinales
\((d_1,\ldots,d_N)\), hojas \((\mu_1,\ldots,\mu_N)\), con
\(\mu_t\in\{\Sigma,\Pi\}\), y los incrementos de su registro
enriquecido. Para \(\lambda\geq0\), definimos
\begin{equation}
 \mathcal A_\lambda(\gamma)
 =\sum_{t=2}^N
 \bigl(1-\delta_{d_t,d_{t-1}}
       +\lambda(1-\delta_{\mu_t,\mu_{t-1}})\bigr).
 \label{vii:eq:accion-ruta-bidireccional}
\end{equation}
El primer sumando cuenta cambios de dirección; el segundo cuenta cambios
de hoja con el peso declarado. Para \(N\leq1\), la suma es cero.
La acción general del sector conserva además sus términos de frontera y
memoria:
\begin{equation}
 \mathcal S[\gamma]
 =\sum_{a\subset\gamma}\mathcal L_a
  +\mathcal B_{\rm in}(\partial_-\gamma)
  +\mathcal B_{\rm out}(\partial_+\gamma)
  +\mathcal M(\gamma).
 \label{vii:eq:accion-enriquecida-frontera}
\end{equation}
Los términos de esta expresión especifican una clase variacional con
extremos tipados y memoria conservada. La dinámica del sector determina
\(\mathcal L_a\), las condiciones admisibles y el término \(\mathcal M\).

Sobre los registros definimos \(\mathcal C_{\rm r}\) como el intercambio
global de hojas; \(\mathcal P_{\rm r}\) como una reflexión de la coordenatización
espacial y de sus direcciones; y \(\mathcal T_{\rm r}\) como la inversión
del orden temporal, la sustitución de cada dirección por su opuesta y el
intercambio de extremos. La reflexión cardinal es una permutación involutiva.
La inversión del registro firmado es
\(C_M(a_1,\ldots,a_N)=(-a_N,\ldots,-a_1)\).
Las hojas, los cocientes y la memoria se transportan con la ruta completa;
la lectura escalar de una palabra constituye una operación posterior.

\begin{proposition}[Invariancia de la acción local de ruta]
\label{vii:prop:invariancia-ruta-bidireccional}
En el dominio de registros anterior,
\[
 \mathcal A_\lambda(\gamma)
 =\mathcal A_\lambda(\mathcal C_{\rm r}\gamma)
 =\mathcal A_\lambda(\mathcal P_{\rm r}\gamma)
 =\mathcal A_\lambda(\mathcal T_{\rm r}\gamma)
 =\mathcal A_\lambda(\mathcal C_{\rm r}\mathcal P_{\rm r}
                              \mathcal T_{\rm r}\gamma).
\]
\end{proposition}
\begin{proof}
Una permutación global del alfabeto de direcciones conserva la igualdad
entre direcciones consecutivas. El intercambio de las dos hojas conserva
la igualdad entre hojas consecutivas. La inversión del orden establece una
biyección entre pares adyacentes, y la sustitución por direcciones opuestas
también conserva sus igualdades. Cada operación conserva los dos recuentos
de \eqref{vii:eq:accion-ruta-bidireccional}; su composición los conserva.
\end{proof}

La proposición determina la simetría del funcional local. Para que una
transformación actúe además sobre un problema dinámico con fronteras fijas,
debe transportar su clase de prolongaciones admisibles y los otros tres
términos de \eqref{vii:eq:accion-enriquecida-frontera}. En particular, dos
rutas con igual representación final pueden conservar distintas clases
variacionales por su memoria.

Una interacción entre dos secciones realizadas en una frontera \(B\)
utiliza un mapa de acoplamiento
\[
 \mathcal I_B:E_B^{(1)}\otimes E_B^{(2)}\longrightarrow E_B^{(12)}
\]
y un término \(\mathcal S_{\rm int}\) de la misma acción. Su covariancia
se expresa por
\[
 \rho_{12}(g)\mathcal I_B
 =\mathcal I_B\bigl(\rho_1(g)\otimes\rho_2(g)\bigr).
\]
Al componer las historias, el registro de la frontera compartida se cuenta
una vez; las memorias de ambos lados se transportan a la sección compuesta.
El acoplamiento y esta condición de covariancia son datos de la interacción
concreta, distintos de las involuciones del registro.

\subsection{Reversión de la evolución realizada}
\label{vii:subsec:reversion-evolucion}

La realización del registro proporciona un espacio de Hilbert real
\(\mathcal H_{\mathbb R}\), una estructura compleja ortogonal
\(J_E\), con \(J_E^2=-I\), y un Hamiltoniano autoadjunto \(H\).
Su dominio \(\mathcal D(H)\) es invariante por \(J_E\); se supone
\(HJ_E=J_EH\) en ese dominio. Sea \(\mathcal C\) una involución
ortogonal que conserva \(\mathcal D(H)\) y satisface
\begin{equation}
 \mathcal C J_E\mathcal C^{-1}=-J_E,\qquad
 \mathcal C H\mathcal C^{-1}=H.
 \label{vii:eq:dominio-reversion-evolucion}
\end{equation}
La constante interna de acción \(\hbar_{\rm int}>0\) conserva su
generación anterior; aquí normaliza el grupo de evolución.

\begin{theorem}[Conjugación temporal del grupo]
\label{vii:thm:reversion-evolucion}
El generador \(A=-J_EH/\hbar_{\rm int}\), con dominio \(\mathcal D(H)\),
es antiadjunto. Su grupo ortogonal, unitario respecto de \(J_E\), cumple
\begin{equation}
 U(t)=\exp(-J_EHt/\hbar_{\rm int}),\qquad
 \mathcal C U(t)\mathcal C^{-1}=U(-t).
 \label{vii:eq:reversion-grupo-bidireccional}
\end{equation}
\end{theorem}
\begin{proof}
La ortogonalidad y \(J_E^2=-I\) implican \(J_E^*=-J_E\).
La invariancia del dominio y la conmutación con \(H\) permiten tratar
\(H\) como operador autoadjunto complejo en el espacio definido por
\(J_E\). Por cálculo funcional, \(-J_EH/\hbar_{\rm int}\) es
antiadjunto y genera el grupo anunciado. En dimensión finita, las mismas
afirmaciones siguen de la serie exponencial y de \(A^*=-A\).
Las hipótesis sobre \(\mathcal C\) dan
\(\mathcal C A\mathcal C^{-1}=-A\). Conjugar el grupo, o su cálculo
funcional, produce \(\exp(-tA)=U(-t)\).
\end{proof}

El mapa que realiza la involución firmada \(C_M\) como \(\mathcal C\)
conserva la orientación temporal y el dominio anterior. La conjugación de
carga de una representación de campos conserva, además, sus propios
caracteres y operadores. La igualdad
\eqref{vii:eq:reversion-grupo-bidireccional} fija la reversión temporal
del grupo bajo las condiciones expresadas.

\subsection{Propagadores del canal cronológico de una ruta}
\label{vii:subsec:green-cronologico}

La inscripción isométrica del registro admite las prolongaciones unitarias
del teorema~\ref{vii:thm:observador-entrelazador}. Fijemos una realización
bilateral de un canal cronológico con fibras \(\mathcal H_n\),
\(n\in\mathbb Z\), y transportes unitarios
\(U_n:\mathcal H_n\to\mathcal H_{n+1}\).
La elección incluye los espacios auxiliares cuando se requieren; conserva
la ruta y sus registros anteriores. Denotamos por \(U(n,k)\) el transporte
de \(k\) a \(n\), con
\[
 U(k,k)=I,\qquad U(n,k)U(k,l)=U(n,l),\qquad
 U(n,k)^*=U(k,n).
\]
Para \(n>k\), es el producto \(U_{n-1}\cdots U_k\); para \(n<k\),
es el inverso del transporte de \(n\) a \(k\).

El espacio de fuentes \(\mathcal J_c\) contiene las secciones de soporte
finito. El espacio de respuestas \(\mathcal F\) contiene todas las
secciones de esas fibras. La acción cuadrática de la ruta, en unidades
internas de esta realización, utiliza
\[
 (D\psi)_n=\psi_{n+1}-U_n\psi_n,\qquad
 E[\psi]=\sum_n\|(D\psi)_n\|^2
\]
para secciones de soporte finito. El operador cinético con la convención
de signo \(L=-D^*D\) es
\begin{equation}
 (L\psi)_n=U_n^*\psi_{n+1}-2\psi_n+U_{n-1}\psi_{n-1}.
 \label{vii:eq:diferencia-covariante-bidireccional}
\end{equation}
En efecto, \((D^*z)_n=z_{n-1}-U_n^*z_n\), y la sustitución produce
esa fórmula. La variación de \(E\) es
\(-2\operatorname{Re}\langle\delta\psi,L\psi\rangle\).

\begin{theorem}[Green exactos de la diferencia covariante]
\label{vii:thm:green-cronologico}
Escribiendo \(x_+=\max(x,0)\), las aplicaciones
\begin{align}
 (G_{\rm ret}j)_n&=\sum_k(n-k)_+U(n,k)j_k,\label{vii:eq:green-ret-cronologico}\\
 (G_{\rm adv}j)_n&=\sum_k(k-n)_+U(n,k)j_k\label{vii:eq:green-adv-cronologico}
\end{align}
de \(\mathcal J_c\) en \(\mathcal F\) satisfacen
\[
 LG_{\rm ret}j=LG_{\rm adv}j=j.
\]
La respuesta retardada se anula antes y en la primera posición de la
fuente; la avanzada se anula después y en la última. Son las únicas
soluciones que se anulan suficientemente lejos en la dirección respectiva.
Para \(\psi\) de soporte finito también se cumple
\(G_{\rm ret}L\psi=G_{\rm adv}L\psi=\psi\).
Respecto del emparejamiento con fuentes de soporte finito,
\(G_{\rm adv}=G_{\rm ret}^{\dagger}\).
\end{theorem}
\begin{proof}
Cada suma tiene tantos términos como posiciones de la fuente.
La composición de los transportes reduce la aplicación de \(L\) a
la identidad escalar
\[
 (n+1-k)_+-2(n-k)_++(n-1-k)_+=\delta_{nk}.
\]
La función \((k-n)_+\) satisface la misma identidad. Esto prueba
las dos ecuaciones y los soportes. Para la unicidad, una solución
homogénea que se anula en dos posiciones consecutivas queda determinada
como cero en la siguiente por \eqref{vii:eq:diferencia-covariante-bidireccional};
la recurrencia, usando transportes invertibles, se prolonga en ambas
direcciones. La diferencia de dos soluciones con el mismo soporte causal
se anula suficientemente lejos y, por tanto, es cero.
Aplicar este argumento a \(\psi\) y a \(G_{\rm ret}L\psi\), o al
caso avanzado, da las identidades por la derecha sobre soporte finito.
Finalmente, para dos fuentes compactas, la suma doble del emparejamiento
es finita y
\[
 \bigl((k-n)_+U(k,n)\bigr)^*=(k-n)_+U(n,k).
\]
Intercambiar los índices demuestra la adjunción indicada.
\end{proof}

Las respuestas pertenecen a \(\mathcal F\); en general crecen
linealmente fuera de la fuente y no pertenecen a \(\ell^2\).
El dominio de soporte finito y su emparejamiento especifican la adjunción
de Green utilizada aquí. El resultado construye la propagación del canal
cronológico desde su acción; una realización de campos espaciales conserva
además su operador, geometría y dominio de propagación propios.

Sea \(C_0\) una involución isométrica, lineal o antilineal, de
\(\mathcal H_0\). El transporte desde la fibra de referencia construye
\[
 C_n=U(-n,0)C_0U(0,n):\mathcal H_n\longrightarrow\mathcal H_{-n},
 \qquad (\mathcal C\psi)_n=C_{-n}\psi_{-n}.
\]
Se tiene \(C_{-n}C_n=I\) y
\(C_{n+1}U_n=U_{-n-1}^*C_n\), por composición.
En particular, \(\mathcal C^2=I\), y el cambio \((n,k)\mapsto(-n,-k)\)
en las sumas anteriores prueba
\begin{equation}
 \mathcal C L=L\mathcal C,\qquad
 \mathcal C G_{\rm ret}\mathcal C^{-1}=G_{\rm adv}.
 \label{vii:eq:intercambio-green-cronologico}
\end{equation}
Cuando \(C_0\) es la realización declarada de la involución del registro,
estas fórmulas transportan esa involución a todo el canal, sin identificar
el retorno de fase con un reinicio de memoria.

\subsection{Intercambio causal y descomposición de la respuesta}
\label{vii:subsec:green-realizacion-campos}

En una realización de campos, sea
\(L:\mathcal D(L)\subset\mathcal H_\Phi\to\mathcal H_J\)
un operador con Green retardado y avanzado únicos en los dominios de
fuentes admisibles. Denotamos por \(J^+(S)\) y \(J^-(S)\) los
dominios causales futuro y pasado de un soporte \(S\). Las condiciones son
\[
 \begin{gathered}
 LG_{\rm ret}=LG_{\rm adv}=I,\\
 \operatorname{supp}(G_{\rm ret}j)\subset J^+(\operatorname{supp}j),\\
 \operatorname{supp}(G_{\rm adv}j)\subset J^-(\operatorname{supp}j).
 \end{gathered}
\]
Sean \(C_\Phi,C_J\) involuciones isométricas, ambas lineales o ambas
antilineales, que transportan esos dominios,
invierten sus orientaciones causales y satisfacen
\(LC_\Phi=C_JL\). Se conserva un emparejamiento de Green para el cual
\(G_{\rm adv}=G_{\rm ret}^{\dagger}\).

\begin{proposition}[Conjugación de los Green y paridades]
\label{vii:prop:green-paridades}
Con las condiciones anteriores,
\[
 C_\Phi G_{\rm ret}C_J^{-1}=G_{\rm adv}.
\]
Sobre el espacio real de operadores de respuesta, definamos
\[
 \Theta(K)=C_\Phi K C_J^{-1},\qquad
 \mathsf P_\pm(K)=\tfrac12(K\pm\Theta(K)).
\]
Son proyecciones complementarias para la involución \(\Theta\), y
\begin{equation}
 G_{\rm sym}=\tfrac12(G_{\rm ret}+G_{\rm adv}),\qquad
 G_{\rm rad}=\tfrac12(G_{\rm ret}-G_{\rm adv})
 \label{vii:eq:descomposicion-green}
\end{equation}
satisfacen \(\Theta(G_{\rm sym})=G_{\rm sym}\),
\(\Theta(G_{\rm rad})=-G_{\rm rad}\),
\(G_{\rm sym}^{\dagger}=G_{\rm sym}\) y
\(LG_{\rm rad}=0\). El soporte de la respuesta simétrica está contenido
en la unión de los dos dominios causales.
\end{proposition}
\begin{proof}
El operador conjugado resuelve la ecuación de fuente, pues
\[
 LC_\Phi G_{\rm ret}C_J^{-1}=C_JLG_{\rm ret}C_J^{-1}=I.
\]
La inversión de la orientación lleva su soporte al dominio avanzado;
su unicidad lo identifica con \(G_{\rm adv}\). Las involuciones dan
\(\Theta^2=I\); expandir \((I\pm\Theta)^2\) y
\((I+\Theta)(I-\Theta)\) demuestra las afirmaciones sobre proyecciones.
La adjunción de Green demuestra la simetría y la diferencia de las dos
ecuaciones de fuente da \(LG_{\rm rad}=0\). La suma de dos respuestas
tiene soporte contenido en la unión de sus soportes.
\end{proof}

Si las involuciones son antilineales, \(\Theta\) es lineal sobre
escalares reales; ése es el espacio operatorio usado en la proposición.
La conjugación del grupo temporal, la condición de soporte y la
adjunción de Green son tres propiedades distintas, transmitidas por los
mapas de realización correspondientes.

La reducción causal de la actualización bilateral, desarrollada en la
sección sobre conservación y recuperación, identifica otro origen preciso
del retardo: la eliminación de una componente de memoria que sigue
participando en la evolución conjunta. Su núcleo discreto conserva el
orden de las interacciones y la contribución de la memoria inicial.
Los operadores de Green aquí construidos resuelven, en cambio, una
ecuación de propagación con condiciones de soporte. Ambas construcciones
pueden componerse cuando se fija el mapa de realización que identifica
sus dominios y operadores; sus respectivas condiciones de causalidad
permanecen explícitas.

\subsection{Proyección de corrientes absorbidas y variación de la acción}
\label{vii:subsec:proyeccion-absorcion}

Sea \(\gamma_\partial:\mathcal H_\Phi\to\mathcal H_\partial\)
la lectura de frontera, distinta de la ruta \(\gamma\). Definimos
\[
 \mathcal A_\partial=\gamma_\partial G_{\rm rad}.
\]
La ventana finita fija el espacio de fuentes y la lectura de frontera;
los Green conservan su dominio de respuestas. En un nivel finito con
productos internos fijados, su inversa de
Moore--Penrose determina
\begin{equation}
 P_{\rm abs}=I-\mathcal A_\partial^+\mathcal A_\partial.
 \label{vii:eq:proyector-absorcion}
\end{equation}

\begin{theorem}[Condición de absorción y respuesta de borde]
\label{vii:thm:proyeccion-absorcion}
\(P_{\rm abs}\) es la proyección ortogonal sobre
\(\ker\mathcal A_\partial\). Para cualquier fuente \(j\), la fuente
\(j_{\rm abs}=P_{\rm abs}j\) cumple
\begin{equation}
 \gamma_\partial G_{\rm ret}j_{\rm abs}
 =\gamma_\partial G_{\rm adv}j_{\rm abs}.
 \label{vii:eq:igualdad-borde-absorbido}
\end{equation}
Es la fuente absorbida más próxima a \(j\) para el producto interno
declarado. Existen fuentes absorbidas no nulas exactamente cuando
\(\operatorname{rank}\mathcal A_\partial<\dim\mathcal H_J\).
\end{theorem}
\begin{proof}
La descomposición ortogonal
\(\mathcal H_J=\ker\mathcal A_\partial\oplus
 \operatorname{im}\mathcal A_\partial^*\) muestra que
\(\mathcal A_\partial^+\mathcal A_\partial\) es la identidad
en el segundo sumando y cero en el primero. Su complemento es el
proyector anunciado. Por definición,
\(2\mathcal A_\partial j_{\rm abs}
 =\gamma_\partial(G_{\rm ret}-G_{\rm adv})j_{\rm abs}=0\).
Para \(k\in\ker\mathcal A_\partial\), la descomposición ortogonal da
\[
 \|j-k\|^2=\|(I-P_{\rm abs})j\|^2+\|P_{\rm abs}j-k\|^2.
\]
El mínimo único corresponde a \(k=P_{\rm abs}j\).
La equivalencia final es el teorema de rango y nulidad.
\end{proof}

En el emparejamiento de fuente y campo, la acción simétrica es
\(\mathscr S[j]=\tfrac12\langle j,G_{\rm sym}j\rangle\).
La adjunción declarada hace real esa cantidad. Para variaciones reales,
\[
 \mathrm d\mathscr S[j](h)
 =\operatorname{Re}\langle h,G_{\rm sym}j\rangle.
\]
La acción restringida a fuentes absorbidas se escribe sobre el espacio
ambiente como
\begin{equation}
 \mathscr S_{\rm abs}[j]
 =\tfrac12\langle P_{\rm abs}j,G_{\rm sym}P_{\rm abs}j\rangle,
 \quad
 \mathrm d\mathscr S_{\rm abs}[j](h)
 =\operatorname{Re}\langle P_{\rm abs}h,G_{\rm sym}P_{\rm abs}j\rangle.
 \label{vii:eq:variacion-restringida-absorcion}
\end{equation}
Cuando el emparejamiento identifica campo y fuente por Riesz, el gradiente
ambiente es \(P_{\rm abs}G_{\rm sym}P_{\rm abs}j\).
La respuesta de campo sigue siendo \(G_{\rm sym}j_{\rm abs}\); el
proyector del gradiente expresa la restricción de las variaciones.
Las fórmulas se demuestran expandiendo el término cuadrático de
\(j+th\) y utilizando la autoadjunción. Un término material
\(\mathscr S_{\rm matter}\) añade su propio diferencial.

\subsection{Condición absorbente explícita en una ventana cronológica}
\label{vii:subsec:absorcion-momentos}

El canal de la subsección~\ref{vii:subsec:green-cronologico} permite
evaluar esa condición sin introducir una respuesta objetivo. Sea una
fuente apoyada en \(0,\ldots,N\), con \(N\geq1\). Transportemos
sus componentes a la fibra de referencia:
\[
 \widehat j_k=U(0,k)j_k,\qquad
 M_0(j)=\sum_{k=0}^N\widehat j_k,\qquad
 M_1(j)=\sum_{k=0}^Nk\widehat j_k.
\]
La identidad \((n-k)_+-(k-n)_+=n-k\) da
\begin{equation}
 (G_{\rm ret}-G_{\rm adv})j\big|_n
 =U(n,0)\bigl(nM_0(j)-M_1(j)\bigr).
 \label{vii:eq:diferencia-green-momentos}
\end{equation}

\begin{proposition}[Absorción y dos momentos transportados]
\label{vii:prop:absorcion-momentos}
Si la traza evalúa el campo en dos posiciones distintas \(a,b\),
la igualdad de trazas avanzada y retardada equivale a
\(M_0(j)=M_1(j)=0\). En una fibra de dimensión \(d\), el espacio
de tales fuentes tiene dimensión \((N-1)d\). Para \(N\geq2\),
contiene las fuentes no nulas con componentes transportadas
\((z,-2z,z,0,\ldots,0)\), \(z\neq0\).
\end{proposition}
\begin{proof}
La invertibilidad de \(U(a,0)\) y \(U(b,0)\) reduce la condición
de borde a \(aM_0-M_1=0\) y \(bM_0-M_1=0\).
Restar las ecuaciones da \((a-b)M_0=0\), y después \(M_1=0\).
El mapa de los dos momentos es sobreyectivo: para valores prescritos
\((z_0,z_1)\), basta tomar
\(\widehat j_0=z_0-z_1\), \(\widehat j_1=z_1\) y las demás
componentes cero. Su rango es \(2d\); la nulidad es
\((N+1)d-2d\). La fuente indicada tiene ambos momentos nulos.
\end{proof}

En el marco transportado, escribamos
\[
 B_N=\begin{pmatrix}1&1&\cdots&1\\0&1&\cdots&N\end{pmatrix}.
\]
El proyector absorbente queda evaluado por
\begin{equation}
 P_N=\bigl[I-B_N^{\mathsf T}(B_NB_N^{\mathsf T})^{-1}B_N\bigr]
       \otimes I_{\mathcal H_0}.
 \label{vii:eq:proyector-momentos-cronologicos}
\end{equation}
Las dos filas de \(B_N\) son independientes para \(N\geq1\), por lo
que la inversa existe. El transporte unitario entre marcos lleva este
proyector a las fibras originales. Para \(N=2\),
\[
 P_2=\frac16
 \begin{pmatrix}1&-2&1\\-2&4&-2\\1&-2&1\end{pmatrix}
 \otimes I_{\mathcal H_0}.
\]
Esta construcción produce una clase absorbida no nula y su operación de
proyección en el canal cronológico. Su amplitud y su identificación con una
corriente material conservan las reglas del sector físico correspondiente.

\subsection{Refinamiento y continuidad de la condición de borde}
\label{vii:subsec:refinamiento-absorcion}

Sean \(R_m^\Phi,R_m^J,R_m^\partial\) los mapas de refinamiento de
campos, fuentes y fronteras. Supongamos que entrelazan los dos Green y la
traza. Por composición se obtiene
\begin{equation}
 R_m^\Phi G_{{\rm sym},m}=G_{{\rm sym},m+1}R_m^J,\qquad
 R_m^\partial\mathcal A_{\partial,m}
 =\mathcal A_{\partial,m+1}R_m^J.
 \label{vii:eq:refinamiento-lectura-absorcion}
\end{equation}
La segunda igualdad lleva una fuente absorbida a otra fuente absorbida.
La naturalidad de los proyectores conserva adicionalmente la estructura
ortogonal de las realizaciones.

\begin{proposition}[Naturalidad de la proyección absorbente]
\label{vii:prop:naturalidad-proyeccion-absorcion}
Para una isometría de fuentes \(R=R_m^J\), se tiene
\(P_{{\rm abs},m+1}R=RP_{{\rm abs},m}\) exactamente cuando
\[
 R(\ker\mathcal A_{\partial,m})
       \subset\ker\mathcal A_{\partial,m+1},\qquad
 R((\ker\mathcal A_{\partial,m})^\perp)
       \subset(\ker\mathcal A_{\partial,m+1})^\perp.
\]
Con esas condiciones en todos los niveles, los proyectores definen una
proyección ortogonal en la completación del límite inductivo isométrico
de los espacios de fuentes.
\end{proposition}
\begin{proof}
Se descompone cada fuente en sus componentes en el núcleo y en su
complemento ortogonal. Las dos inclusiones implican que aplicar la
proyección antes o después del refinamiento conserva respectivamente la
primera componente y anula la segunda. Recíprocamente, la igualdad de
operadores aplicada a cada componente da las dos inclusiones.
En el límite inductivo, la igualdad permite definir un operador sobre la
unión de los niveles sin depender del representante. Todos los proyectores
tienen norma a lo sumo uno, de modo que el operador se extiende a la
completación. Sus identidades de autoadjunción e idempotencia, válidas en
la unión densa, pasan por continuidad a la completación.
\end{proof}

Para prolongar también una respuesta de Green a una completación se
conservan sus dominios y las estimaciones de continuidad correspondientes.
La propagación cronológica de soporte finito y la factorización de la
respuesta por la imagen de borde, demostrada en
\ref{vii:thm:observador-factorizacion}, mantienen así sus dominios
respectivos durante el refinamiento.

La realización física de emisión y absorción compone estos operadores con
las corrientes materiales, la acción y las condiciones causales del sector.
El entrelazamiento debe conservar la involución de ruta, el emparejamiento
de fuente y campo, la acción transportada y las condiciones de frontera a
cada nivel. La construcción precedente aporta la reversión del grupo, los
Green del canal cronológico, la descomposición temporal y una proyección
absorbente explícita; las realizaciones espaciales emplean sus mapas y
condiciones de causalidad propios.
